\chapter{Literature Review}
\label{ch:literature}

\begin{chapterabstract}
This chapter surveys the research streams on which the thesis builds:
vision-language modelling, efficient and edge inference, temporal redundancy
reduction in video, real-time object detection and multi-object tracking,
semantic representations and concept drift, and adaptive conditional
computation. Each section closes by identifying the gap that motivates a
\emph{semantic-aware} scheduling layer, and the chapter ends with a
comparative positioning of the proposed approach.
\end{chapterabstract}

\section{Vision-Language Models}
\label{sec:lit-vlm}

Vision-language models learn a joint representation of images and text that
supports open-vocabulary recognition and reasoning. CLIP~\cite{Radford2021CLIP}
established the contrastive pre-training paradigm, aligning an image encoder
and a text encoder in a shared embedding space using large-scale web data; its
normalised embeddings have since become a de-facto semantic descriptor for
images. Generative VLMs extend this paradigm to free-form language output:
BLIP-2~\cite{Li2023BLIP2} bridges a frozen vision encoder to a frozen large
language model through a lightweight querying transformer, while
LLaVA~\cite{Liu2023LLaVA} fine-tunes a projection between visual features and
an instruction-tuned language model. These architectures inherit the
transformer backbone~\cite{Vaswani2017Attention,Dosovitskiy2021ViT} and, with
it, a per-inference cost that scales with input resolution and sequence
length.

For edge deployment, compact VLMs are essential. Moondream~\cite{Moondream2024}
is a representative tiny VLM (on the order of a few billion parameters)
designed for captioning and visual question answering on modest hardware. Even
so, a single Moondream invocation remains orders of magnitude more expensive
than the lightweight perception primitives discussed below, which is precisely
why \emph{when} to invoke it matters.

\section{Vision-Language Models in Robotics}
\label{sec:lit-robotics}

A growing body of work embeds VLMs into robotic control. SayCan~\cite{Ahn2022SayCan}
grounds language-model plans in learned robotic affordances;
RT-2~\cite{Brohan2023RT2} and OpenVLA~\cite{Kim2024OpenVLA} train
vision-language-action models that map observations and instructions directly
to control. Closer to the temporal setting of this thesis, agentic systems
increasingly maintain a memory of the scene rather than reasoning from each
observation in isolation: VideoAgent~\cite{Wang2024VideoAgent} couples a
multimodal agent with an explicit memory to answer queries over long videos,
and VideoARM~\cite{CognitiveNav} performs agentic reasoning over a
hierarchical memory for long-form video understanding. These systems
demonstrate the value of multimodal reasoning and persistent semantic memory
for understanding continuous video, but they typically assume access to
substantial compute and treat perception as a per-step query. None incorporates an explicit notion of \emph{semantic
stability} that would let the agent skip reasoning when its understanding of
the scene is unchanged --- the omission this thesis targets.

\section{Efficient and Edge Inference}
\label{sec:lit-efficient}

Reducing the cost of neural inference is a mature field. Architectural
efficiency is pursued through compact backbones such as
MobileNets~\cite{Howard2017MobileNets}. Post-training compression includes
pruning and quantization~\cite{Han2016DeepCompression} and low-precision
matrix multiplication tailored to transformers~\cite{Dettmers2022LLMint8}.
Knowledge distillation~\cite{Hinton2015Distillation} transfers the behaviour
of a large model into a smaller one; TinyCLIP~\cite{Wu2023TinyCLIP} applies
this idea to the embedding model itself, distilling a compact CLIP suitable
for edge use --- directly relevant to the lightweight embedding the monitor
relies on. The broader case for moving inference to the network edge, and the
energy, latency, and bandwidth constraints that motivate it, is set out in
surveys of edge intelligence~\cite{Zhou2019EdgeIntelligence}. These techniques
reduce the cost of an \emph{individual} forward pass, and are complementary to
this thesis: they make each VLM call cheaper, whereas the proposed scheduler
reduces the \emph{number} of calls. The two axes multiply, and a complete edge
system would exploit both.

\section{Temporal Redundancy Reduction in Video}
\label{sec:lit-temporal}

Because video is temporally redundant, classical recognition pipelines first
established that sparse temporal sampling suffices for many tasks: temporal
segment networks~\cite{Wang2016TSN} recognise actions from a handful of
sparsely sampled snippets, and attention-based video
transformers~\cite{Bertasius2021TimeSformer} model temporal structure across
frames. Building on this, several methods learn \emph{which} frames deserve
expensive processing: AdaFrame~\cite{Wu2019AdaFrame} trains a policy to select
informative frames for action recognition; SCSampler~\cite{Korbar2019SCSampler}
samples salient clips using a cheap proxy; skip-convolutions~\cite{Habibian2021Skip}
avoid recomputing convolutional features in regions that do not change between
frames. These approaches operate at the level of raw appearance or low-level
features and are usually trained end-to-end for a fixed downstream task. The
present work differs in two respects: its trigger is a \emph{semantic} drift
signal rather than appearance change, and it is \emph{training-free} at the
policy level, relying on interpretable thresholds rather than a learned
sampler, which suits the heterogeneous, label-scarce conditions of robotic
deployment.

\section{Object Detection and Multi-Object Tracking}
\label{sec:lit-detection}

Real-time object detection has been driven by single-stage detectors since
YOLO~\cite{Redmon2016YOLO}; the YOLOv8 family~\cite{Jocher2023YOLOv8} offers
nano-scale variants that run comfortably on edge accelerators while reporting
object classes and bounding boxes trained on COCO~\cite{Lin2014COCO}.
Multi-object tracking associates detections across frames to maintain stable
identities: SORT~\cite{Bewley2016SORT} couples a Kalman filter with
Hungarian assignment, and ByteTrack~\cite{Zhang2022ByteTrack} improves
robustness by associating low-confidence boxes as well. In this thesis these
components are not contributions in themselves; rather, their outputs ---
object classes, counts, and identity continuity --- are repurposed as cheap,
informative features for estimating semantic drift.

\section{Semantic Representations and Concept Drift}
\label{sec:lit-drift}

Treating the evolving content of a video stream as a non-stationary signal
connects this work to the literature on \emph{concept drift}, surveyed
comprehensively by Gama et al.~\cite{Gama2014Drift}. There, drift detection
governs when a learned model should be updated; here, an analogous drift
signal governs when expensive reasoning should be re-invoked. The semantic
descriptor we monitor is the CLIP embedding~\cite{Radford2021CLIP}, whose
cosine geometry provides a natural, bounded distance between scene states.
The novelty of the present formulation is to \emph{fuse} embedding drift with
symbolic object-set and tracking-identity drift into one decision signal,
rather than relying on a single modality.

\section{Adaptive and Conditional Computation}
\label{sec:lit-adaptive}

Conditional computation allocates network capacity per input. Early-exit
networks such as BranchyNet~\cite{Teerapittayanon2016BranchyNet} terminate
inference once an intermediate classifier is confident; adaptive computation
time~\cite{Graves2017Adaptive} learns how many steps to spend per input; and
dynamic routing methods such as SkipNet~\cite{Wang2018SkipNet} skip layers
conditionally. The survey of Han et al.~\cite{Han2021DynamicSurvey} organises
these into sample-wise, spatial-wise, and temporal-wise dynamics. The
proposed scheduler is a temporal-wise dynamic mechanism operating at the
granularity of \emph{whole-model invocation across frames}, controlled by an
external semantic monitor rather than by internal model confidence --- a design
that keeps the expensive VLM a black box and is therefore agnostic to the
specific model used.

\section{Efficient Video Analytics Systems}
\label{sec:lit-analytics}

A systems-oriented line of research reduces the cost of querying neural
networks over video streams. NoScope~\cite{Kang2017NoScope} accelerates
queries by learning cheap, query-specific cascade models and difference
detectors that filter out frames before invoking an expensive reference
network. Reducto~\cite{Li2020Reducto} pushes filtering onto the camera itself,
discarding frames whose low-level features change little so that only
informative frames reach the analytics backend. These systems established the
template the present work follows --- a cheap filter guarding an expensive
model --- but their filtering criterion is a low-level feature difference tuned
to a fixed query. This thesis replaces that criterion with a \emph{semantic}
drift signal fused from embedding, object, and tracking cues, and guards a
\emph{generative} vision-language model rather than a fixed-output classifier,
so the system reduces semantic reasoning redundancy rather than frame
redundancy for a single predetermined query.

\section{Streaming and Adaptive Multimodal Inference}
\label{sec:lit-streaming}

Most directly related are recent systems that bring adaptivity to multimodal
inference over continuous input. Streaming multimodal
methods~\cite{StreamingVLM} reason over video as it arrives rather than over
pre-segmented clips, raising the question of when continuous reasoning is
actually necessary. AdaVFM~\cite{AdaVFM} makes vision foundation-model
inference adaptive to input difficulty and resource budgets, routing inputs to
cheaper or more expensive computation. QueST~\cite{QueST} frames perception in
terms of representation-level semantic monitoring, providing terminology and
motivation for treating semantic state --- rather than raw frames --- as the
quantity to monitor. The present work shares these systems' goal of
budget-aware multimodal inference but contributes a specific, interpretable,
training-free mechanism: a fused semantic-drift signal mapped through two
thresholds to a three-tier compute policy, with the reasoning model held as a
replaceable black box. Where AdaVFM adapts \emph{within} a model and streaming
methods adapt \emph{how} reasoning proceeds, this thesis decides \emph{whether}
to reason at all, at the granularity of whole-model invocation across frames.

\section{Comparative Positioning}
\label{sec:lit-comparison}

Table~\ref{tab:lit-comparison} contrasts the families above along the axes
most relevant to edge robotic VLM deployment. The proposed approach is the
only one whose decision criterion is an explicit, multi-signal \emph{semantic}
drift estimate, that is training-free, and that treats the reasoning model as
a replaceable black box.

\begin{table}[htbp]
\centering
\caption[Comparison of efficiency approaches]{Positioning of the proposed
semantic-aware scheduler against representative efficiency approaches.}
\label{tab:lit-comparison}
\small
\begin{tabular}{@{}lcccc@{}}
\toprule
\textbf{Approach} & \textbf{Redundancy} & \textbf{Decision} & \textbf{Training} & \textbf{Model-} \\
 & \textbf{targeted} & \textbf{signal} & \textbf{-free} & \textbf{agnostic} \\
\midrule
Compression / quant.\          & Parameter & ---            & No  & No  \\
Frame sampling                 & Frame     & Appearance     & No  & Partial \\
Skip-convolutions              & Pixel     & Pixel change   & No  & No  \\
Early exit                     & Sample    & Model conf.\   & No  & No  \\
NoScope / Reducto              & Frame     & Feature diff.\ & No  & Partial \\
Adaptive VFM (AdaVFM)          & Parameter & Input/budget   & No  & No  \\
Embedding threshold            & Semantic  & Embedding only & Yes & Yes \\
\textbf{Proposed}              & \textbf{Semantic reasoning} & \textbf{Fused drift} & \textbf{Yes} & \textbf{Yes} \\
\bottomrule
\end{tabular}
\end{table}

In summary, prior work makes each inference cheaper or skips low-level
redundancy, but does not reason about \emph{semantic} redundancy at the level
of whole-model invocation using a fused, interpretable signal. This gap
defines the methodology developed next.
